\documentclass[aspectratio=169]{beamer}
\usetheme{Madrid}
\usecolortheme{beaver}

% Packages
\usepackage{graphicx}
\usepackage{hyperref}
\usepackage{amsmath}
\usepackage{minted} % For code highlighting
\usepackage{pgfpages}
% \setbeameroption{show notes on second screen} % Enable for speaker notes

\title[DreamerV4]{Training Agents Inside of Scalable World Models}
\subtitle{Paper Presentation and Practical Code Analysis}
\author{Marco \and Partner}
\institute{Reinforcement Learning Course}
\date{\today}

%% CLAUDE INSTRUCTIONS: 
%% When converting this document to a PowerPoint script, please ensure you allocate roughly 15 minutes for the theoretical Methodology (Slides 1-5) and 15 minutes for the Practical Code Analysis and Demo (Slides 6-10).
%% Please format your output as a slide-by-slide script containing both the bullet points for the screen and the exact speaker notes.

\begin{document}

% Title Slide
\begin{frame}
    \titlepage
\end{frame}

%---------------------------------------------------------
\section{Methodology}
%---------------------------------------------------------

\begin{frame}{The World Model Paradigm}
    \begin{itemize}
        \item \textbf{The Bottleneck:} Standard RL algorithms require millions of real-world interactions, which is slow, expensive, and unsafe for physical robots.
        \item \textbf{The Solution:} Train a "World Model" to simulate the environment. The agent then learns optimal policies entirely in its own "imagination".
        \item \textbf{DreamerV4 Innovation:} Replaces older Recurrent State Space Models (RSSM) with massively scalable \textbf{Block-Causal Transformers}, matching the architecture behind modern LLMs.
    \end{itemize}
\end{frame}

\begin{frame}{Architecture: The Tokenizer}
    \begin{itemize}
        \item \textbf{Objective:} Compress high-dimensional pixels into a compact sequence of latent vectors.
        \item \textbf{Mechanism:} Uses \textbf{Masked Autoencoders (MAE)}. By randomly hiding up to 90\% of the image patches during training, the encoder is forced to learn robust semantic features rather than just compressing pixels.
    \end{itemize}
    \note{Claude: Emphasize why MAE is crucial. Without it, the model might just memorize the environment instead of understanding it.}
\end{frame}

\begin{frame}{Architecture: The Dynamics Model}
    \begin{itemize}
        \item \textbf{Objective:} Predict the future latent states given the past states and the agent's continuous actions.
        \item \textbf{Shortcut Forcing:} Instead of predicting frame-by-frame ($t+1, t+2...$), DreamerV4 uses a denoising schedule ($\tau$) to directly predict future states at $t+k$.
        \item \textbf{Benefit:} Drastically speeds up long-horizon imagination, allowing the agent to plan further ahead without compounding sequential errors.
    \end{itemize}
\end{frame}

%---------------------------------------------------------
\section{Practical Code Analysis}
%---------------------------------------------------------

\begin{frame}{Analyzing the PyTorch Implementation}
    \begin{itemize}
        \item We analyzed an open-source PyTorch implementation applied to \textbf{DMControl} (continuous action spaces), diverging from the paper's original JAX/Minecraft focus.
        \item \textbf{Hardware Reality:} Training this model from scratch requires $\sim$350GB of disk space and 8x High-end GPUs running for 48 hours. 
        \item \textbf{Our Approach:} Focus on the inference architecture and run a live interactive web-server using pre-trained weights.
    \end{itemize}
\end{frame}

\begin{frame}[fragile]{Code Walkthrough: Inside the Encoder}
    The code explicitly separates spatial mixing and temporal mixing inside the \texttt{BlockCausalTransformer}.
    \begin{minted}[fontsize=\footnotesize, bgcolor=lightgray]{python}
class Encoder(nn.Module):
    def forward(self, patch_tokens):
        # 1. Apply Masked Auto-Encoding (MAE) randomly
        proj_masked, mae_mask, keep_prob = self.mae(patch_tokens)
        
        # 2. Process through BlockCausal Layers
        enc = self.transformer(tokens)
        
        # 3. Tanh bottleneck strictly bounds latents in [-1, 1]
        z = torch.tanh(self.bottleneck_proj(enc[:, :, :self.n_latents, :]))
        return z
    \end{minted}
\end{frame}

\begin{frame}{Real-World Engineering Challenges}
    \begin{itemize}
        \item \textbf{Dataset Mismatch:} The repository's preprocessing script (\texttt{preprocess\_dataset.py}) expects \texttt{.png} files, but the HuggingFace dataset only hosts serialized \texttt{.pt} tensor dictionaries.
        \item \textbf{Our Solution:} To run the live demo without crashing due to a \texttt{FileNotFoundError}, we manually patched the Python code to inject a "dummy" black frame when initializing the sequence. 
        \item \textbf{Result:} The dynamics model successfully takes over and "hallucinates" the environment from a blank state based purely on keyboard inputs.
    \end{itemize}
    \note{Claude: Highlight this slide as a demonstration of critical thinking and engineering problem-solving. We didn't just run the code; we debugged and patched it to work under hardware and data constraints.}
\end{frame}

%---------------------------------------------------------
\section{Live Demonstration}
%---------------------------------------------------------

\begin{frame}{Live Demonstration}
    \begin{center}
        \Large \textbf{1. Interactive Jupyter Notebook} \\
        \normalsize Walking through the tensor shapes and architecture in \texttt{dreamer4\_tutorial.ipynb}.
        
        \vspace{0.8cm}
        
        \Large \textbf{2. Real-Time Web Server Inference} \\
        \normalsize Launching \texttt{interactive.py} to control the agent and watch the World Model hallucinate live frames.
    \end{center}
\end{frame}

\end{document}
